% !TeX program = pdflatex
% Documento principal. Compilar desde monografia/ con latexmk -pdf main.tex
% (o mediante LaTeX Workshop). Las referencias utilizan BibTeX.
% No se ha compilado este archivo: la generación de PDF queda a cargo
% del usuario hasta que autorice expresamente al agente a realizarla.

\documentclass[12pt,a4paper]{article}

\usepackage[T1]{fontenc}
\usepackage[utf8]{inputenc}
\usepackage[spanish]{babel}
\usepackage[margin=2.5cm]{geometry}
\usepackage{amsmath,amssymb}
\usepackage{tikz}
\usepackage[hidelinks]{hyperref}
\usepackage{url}

\begin{document}

\begin{titlepage}
  \centering
  {\large Universidad Católica de Salta\par}
  \medskip
  {\large Facultad de Ingeniería e Informática\par}
  {\large Ingeniería en Informática\par}
  \vfill
  {\large Trabajo Final Integrador\par}
  {\large Monografía de investigación\par}
  \bigskip
  {\LARGE\bfseries Influencia de la variabilidad de los tiempos de servicio
    en los sistemas de colas\par}
  \bigskip
  {\Large Análisis comparativo y demostración mediante simulación\par}
  \vfill
  Integrantes: pendientes de completar\par
  \medskip
  Docente: pendiente de completar\par
  \vfill
  Modelos y Simulación --- 2026\par
\end{titlepage}

\tableofcontents
\clearpage

% Incorporar aquí las siguientes secciones a medida que se redacten.
% El resumen se escribirá al finalizar el cuerpo del trabajo.
% !TeX root = ../main.tex

\section{Introducción}
\label{sec:introduccion}

\subsection{Contexto y relevancia del problema}
\label{subsec:contexto}

Los sistemas de colas permiten estudiar situaciones en las que personas, solicitudes o trabajos requieren la atención de recursos con capacidad limitada. Estas situaciones aparecen tanto en servicios de atención al público como en sistemas informáticos y redes de comunicaciones. Aunque sus componentes físicos son diferentes, comparten una misma estructura: las solicitudes llegan, esperan cuando el recurso está ocupado, reciben servicio y abandonan el sistema. La teoría de colas proporciona herramientas para analizar esta estructura y relacionar las características de la demanda y del servicio con indicadores de desempeño, como la cantidad de solicitudes acumuladas y los tiempos de espera \cite{lauberbonomo2021resetting}.

Una primera aproximación al desempeño de estos sistemas consiste en comparar la tasa de llegada con la capacidad media de atención. Sin embargo, esta comparación no describe por sí sola la experiencia de quienes esperan. En una cola con un servidor, la espera también depende de cómo se distribuyen las duraciones de los servicios. Para el modelo M/G/1 con atención por orden de llegada, la fórmula de Pollaczek--Khinchine permite expresar la espera media en función de la tasa de llegada y de los dos primeros momentos del tiempo de servicio. Por lo tanto, además de su media, resulta relevante su variabilidad \cite{virtamoMG1}.

\subsection{Planteamiento y justificación}
\label{subsec:planteamiento}

El problema abordado en esta monografía surge de una limitación de los análisis basados exclusivamente en promedios: dos sistemas pueden tener la misma tasa de llegada y el mismo tiempo medio de atención, pero producir tiempos medios de espera diferentes. En uno de ellos, las duraciones de los servicios pueden ser relativamente uniformes; en el otro, pueden alternarse atenciones breves y prolongadas. Bajo los supuestos del modelo M/G/1 y manteniendo constantes las demás condiciones, una mayor varianza del tiempo de servicio implica una mayor espera media en cola \cite{virtamoMG1}. Esta relación muestra que una descripción de la capacidad basada únicamente en el promedio puede ocultar diferencias relevantes de funcionamiento.

Estudiar esta relación permite integrar contenidos de probabilidad, generación de variables aleatorias, sistemas de colas y simulación de eventos discretos. Asimismo, ofrece una oportunidad para comparar predicciones analíticas con una demostración computacional acotada. En este contexto, la simulación no se plantea como sustituto de una solución teórica disponible, sino como un recurso para representar la evolución del sistema e ilustrar el efecto de modificar la distribución de los tiempos de servicio. El interés del trabajo reside en comprender y explicar un resultado establecido, así como sus condiciones de aplicación, y no en proponer una relación teórica nueva.

A partir de este planteamiento, se formula la siguiente pregunta de investigación:

\begin{quote}
\emph{¿Cómo afecta la variabilidad de los tiempos de servicio al tiempo medio de espera cuando se mantienen constantes la tasa de llegada y el tiempo medio de atención?}
\end{quote}

\subsection{Alcance y enfoque del trabajo}
\label{subsec:alcance}

El análisis se delimita a sistemas con un único servidor, llegadas que siguen un proceso de Poisson de tasa constante y atención según el orden de llegada, conocida como FCFS por la expresión inglesa \emph{first-come, first-served}. Se consideran tiempos de servicio independientes e idénticamente distribuidos, independientes del proceso de llegadas, y una capacidad de espera ilimitada. No se incorporan prioridades, abandonos, rechazos, interrupciones del servicio ni variaciones de la capacidad de atención. Estas simplificaciones permiten concentrar la comparación en la variabilidad del servicio sin introducir simultáneamente otros mecanismos que puedan modificar las esperas.

Dentro de este alcance, se estudia el modelo M/G/1 y se utilizan M/D/1 y M/M/1 como casos particulares, correspondientes a tiempos de servicio constantes y exponenciales, respectivamente \cite{virtamoMG1}. La comparación mantiene fija la tasa de llegada y la media de servicio. Se trabaja bajo condiciones de estabilidad, con una utilización del servidor menor que uno, y con un segundo momento del tiempo de servicio finito, de modo que resulte aplicable el análisis de la espera media. El indicador principal es el tiempo de espera en cola, es decir, el período transcurrido desde la llegada de una solicitud hasta el inicio de su atención; se distingue del tiempo total de permanencia, que incluye también el servicio.

La monografía adopta un enfoque de revisión bibliográfica y análisis comparativo, complementado con una demostración mediante simulación de eventos discretos cuyo diseño se precisará en la metodología. La revisión permite fundamentar los modelos y discutir sus supuestos; la demostración se orienta a ilustrar la relación estudiada mediante parámetros hipotéticos explícitos. Por este motivo, no se pretende validar un servicio real ni formular recomendaciones operativas para una organización particular. Las conclusiones se limitarán a la evidencia analizada y a las condiciones del modelo, evitando extender el resultado a cualquier disciplina de atención o a cualquier medida de desempeño.

% !TeX root = ../main.tex

\section{Objetivos}
\label{sec:objetivos}

\subsection{Objetivo general}
\label{subsec:objetivo-general}

Analizar la influencia de la variabilidad de los tiempos de servicio en el tiempo medio de espera en cola de un sistema M/G/1 con disciplina FCFS, manteniendo constantes la tasa de llegada y el tiempo medio de atención, mediante una revisión bibliográfica, una comparación analítica y una demostración con simulación de eventos discretos.

\subsection{Objetivos específicos}
\label{subsec:objetivos-especificos}

\begin{enumerate}
  \item Describir los fundamentos de los sistemas de colas con un servidor, identificando los supuestos de los modelos estudiados y distinguiendo la media y la variabilidad del tiempo de servicio, la utilización del servidor y el tiempo de espera en cola.

  \item Explicar, a partir de la fórmula de Pollaczek--Khinchine, la relación entre la variabilidad del tiempo de servicio y la espera media, y comparar los casos M/D/1 y M/M/1 como particularizaciones del modelo M/G/1 bajo iguales condiciones de llegada y media de servicio.

  \item Comparar críticamente los aportes de la bibliografía seleccionada sobre variabilidad y desempeño de sistemas de colas, diferenciando los resultados según sus supuestos, disciplinas de atención y medidas de desempeño.

  \item Desarrollar una demostración reproducible mediante simulación de eventos discretos que compare distribuciones de servicio con igual media y distinta variabilidad, y contrastar las estimaciones del tiempo medio de espera con los valores analíticos de referencia, considerando la incertidumbre estadística de las estimaciones.

  \item Interpretar los resultados de la comparación e identificar sus implicaciones y limitaciones, delimitando las condiciones bajo las cuales las conclusiones pueden aplicarse a los sistemas estudiados.
\end{enumerate}

% !TeX root = ../main.tex

\section{Marco conceptual}
\label{sec:marco-conceptual}

\subsection{Componentes y notación del sistema}
\label{subsec:componentes-cola}

Un sistema de colas se describe mediante las entidades que requieren atención, el proceso de llegadas, el espacio de espera y los recursos que prestan el servicio. En este trabajo, una entidad representa una solicitud o trabajo, sin asociarla a una organización concreta. La llegada marca su ingreso al sistema; si el servidor está libre, comienza la atención inmediatamente. En caso contrario, la entidad permanece en la cola hasta que puede ser atendida. La salida ocurre al finalizar el servicio. Por tanto, la cola comprende únicamente las entidades que esperan, mientras que el sistema incluye también a la que está recibiendo atención.

La figura~\ref{fig:cola-un-servidor} separa estos componentes para evitar confundir acumulación y atención. El bloque de espera representa una capacidad ilimitada, no una cantidad fija de lugares. La secuencia muestra el recorrido lógico: pasar por la cola no implica una espera positiva cuando el servidor está disponible.

\begin{figure}[htbp]
  \centering
  % Fragmento editable. Requiere tikz; no requiere bibliotecas adicionales.
% Esquema de componentes: no representa una realización ni un número de clientes.
% Ancho nominal: 0.96\textwidth, incluidos los bloques extremos.
\begin{tikzpicture}[
  x=\textwidth,y=1cm,font=\small,
  bloque/.style={draw=black!75,fill=white,line width=0.5pt,
    align=center,inner sep=0pt,minimum height=1.35cm},
  flujo/.style={->,draw=black!75,line width=0.6pt}
]
  \node[bloque,text width=0.18\textwidth] (llegadas) at (0.09,0)
    {Llegadas\\Poisson\\$\lambda$};
  \node[bloque,text width=0.24\textwidth] (cola) at (0.37,0)
    {Cola FCFS\\Espera $W_q$};
  \node[bloque,text width=0.20\textwidth] (servidor) at (0.67,0)
    {Un servidor\\Servicio $S$};
  \node[bloque,text width=0.12\textwidth] (salida) at (0.90,0)
    {Salida};
  \draw[flujo] (llegadas.east) -- (cola.west);
  \draw[flujo] (cola.east) -- (servidor.west);
  \draw[flujo] (servidor.east) -- (salida.west);
  \node[align=center] at (0.37,-1.15) {Capacidad de espera ilimitada};
  \draw[draw=black!60,line width=0.45pt]
    (0.25,-1.65) -- (0.25,-1.85) -- (0.77,-1.85) -- (0.77,-1.65);
  \node[align=center,anchor=north] at (0.51,-1.95)
    {Permanencia en el sistema: $T=W_q+S$};
\end{tikzpicture}

  \caption{Esquema conceptual de una cola M/G/1 con disciplina FCFS. Elaboración propia basada en la descripción del modelo de Virtamo \cite{virtamoMG1}. Los bloques representan componentes, no cantidades observadas de entidades.}
  \label{fig:cola-un-servidor}
\end{figure}

La notación de Kendall, en su forma abreviada $A/S/c$, resume el proceso de llegadas, la distribución de los tiempos de servicio y el número de servidores. En esta expresión, $S$ ocupa una posición descriptiva; más adelante se utilizará también como símbolo de la variable aleatoria de servicio. La letra M identifica el caso sin memoria: llegadas Poisson en la primera posición y servicio exponencial en la segunda. D indica una duración determinista y G una distribución general. Así, M/G/1 describe llegadas Poisson, servicio general y un solo servidor \cite{virtamoMG1}. La abreviatura no reemplaza la especificación de la capacidad de espera ni de la disciplina de atención, que deben declararse por separado.

En particular, M/D/1 y M/M/1 son casos de la familia M/G/1: el primero fija una duración constante del servicio y el segundo adopta una distribución exponencial. La letra G no designa una tercera distribución elegida para comparar, sino que permite analizar distintas distribuciones mediante un marco común. Una distribución adicional para la futura demostración requerirá una elección y una parametrización explícitas en la metodología.

\subsection{Llegadas, servicio y disciplina de atención}
\label{subsec:llegadas-servicio-fcfs}

Se supone un proceso de Poisson homogéneo de tasa $\lambda>0$. Esto significa que el número de llegadas en un intervalo de duración $t$ tiene distribución Poisson de media $\lambda t$ y que los conteos en intervalos disjuntos son independientes. Los tiempos entre llegadas consecutivas son independientes y tienen distribución exponencial de parámetro $\lambda$, con media $1/\lambda$. La tasa se expresa en entidades por unidad de tiempo. Su constancia describe la intensidad media de la demanda, pero no implica que las llegadas estén igualmente espaciadas.

El tiempo de servicio se representa mediante una variable aleatoria no negativa $S$. Las duraciones de las sucesivas atenciones son independientes e idénticamente distribuidas e independientes del proceso de llegadas. Se adopta FCFS, del inglés \emph{first-come, first-served}: las entidades comienzan a ser atendidas en el orden en que llegan. El servicio es no expropiativo, es decir, una atención iniciada no se interrumpe para atender otra solicitud. Además, el servidor trabaja siempre que hay entidades pendientes; no se incluyen pausas, vacaciones ni tiempos adicionales de preparación.

Se considera una fuente potencial infinita, por lo que la intensidad de llegada no disminuye por la cantidad de entidades dentro del sistema. La capacidad de espera es ilimitada y no hay rechazos ni abandonos: toda llegada permanece hasta completar su atención. Estos supuestos permiten usar la misma tasa $\lambda$ para los ingresos y, en régimen estacionario, para las salidas. El marco corresponde a una cola abierta aislada, no a una red cerrada ni a un modelo con prioridades.

\subsection{Media, variabilidad y utilización}
\label{subsec:variabilidad-utilizacion}

La media del servicio y su tasa media equivalente se definen como
\begin{equation}
  m=\mathrm{E}[S], \qquad \mu=\frac{1}{m}, \qquad 0<m<\infty.
  \label{eq:media-tasa-servicio}
\end{equation}
Si $m$ se mide en minutos por entidad, $\mu$ se expresa en entidades por minuto. Definir $\mu$ como el inverso de la media no presupone que el servicio sea exponencial. La dispersión de las duraciones se caracteriza mediante
\begin{equation}
  \operatorname{Var}(S)=\mathrm{E}[(S-m)^2],
  \qquad C_s^2=\frac{\operatorname{Var}(S)}{m^2}.
  \label{eq:variabilidad-servicio}
\end{equation}
La varianza tiene unidades de tiempo al cuadrado. El coeficiente de variación $C_s=\sqrt{\operatorname{Var}(S)}/m$ y su cuadrado $C_s^2$ son adimensionales; este último permite expresar la dispersión relativa sin depender de la unidad temporal. El segundo momento satisface
\begin{equation}
  \mathrm{E}[S^2]=\operatorname{Var}(S)+m^2=m^2(1+C_s^2).
  \label{eq:segundo-momento-servicio}
\end{equation}
Estas relaciones permiten separar el nivel medio de atención de su variabilidad \cite{virtamoMG1}.

La intensidad de tráfico es
\begin{equation}
  \rho=\lambda m=\frac{\lambda}{\mu}<1.
  \label{eq:utilizacion-estabilidad}
\end{equation}
Es una cantidad adimensional y, bajo los supuestos adoptados en régimen estacionario, representa la fracción de tiempo que el servidor permanece ocupado. La desigualdad establece la condición de estabilidad: el trabajo que ingresa por unidad de tiempo debe ser inferior a la capacidad disponible \cite{virtamoMG1}. No significa que nunca se formen colas; aun con capacidad media suficiente, la irregularidad de llegadas y servicios produce esperas. Para estudiar una espera media finita se exige además $\mathrm{E}[S^2]<\infty$. Estabilidad y finitud de ese promedio son requisitos distintos: una media de servicio finita no garantiza un segundo momento finito.

\subsection{Medidas de desempeño y ley de Little}
\label{subsec:medidas-little}

Se denota por $W_q$ el tiempo aleatorio que una entidad espera antes de iniciar su atención, y por $T$ su tiempo total en el sistema. Para esa misma entidad,
\begin{equation}
  T=W_q+S, \qquad \mathrm{E}[T]=\mathrm{E}[W_q]+m.
  \label{eq:tiempo-sistema}
\end{equation}
El subíndice $q$ identifica la espera en cola. Esta distinción evita atribuir al servicio propio una demora que corresponde a trabajos anteriores. Una entidad atendida inmediatamente tiene $W_q=0$, aunque su permanencia total incluya el servicio.

Sean $N_q$ y $N$ las cantidades aleatorias de entidades en espera y en el sistema, respectivamente, observadas en régimen estacionario. Sus medias se escriben $L_q=\mathrm{E}[N_q]$ y $L=\mathrm{E}[N]$. La ley de Little relaciona estas cantidades con los tiempos medios:
\begin{equation}
  L_q=\lambda\mathrm{E}[W_q], \qquad L=\lambda\mathrm{E}[T].
  \label{eq:ley-little}
\end{equation}
Se aplica aquí en condiciones estacionarias, con tasas de flujo coherentes y promedios finitos \cite{virtamoMG1}. No es una igualdad entre observaciones individuales ni una garantía de coincidencia exacta en una ejecución finita. Tampoco requiere por sí misma distribuciones Poisson o exponenciales: su generalidad es mayor que la del modelo empleado para calcular la espera. En este sistema, las ecuaciones anteriores implican $L=L_q+\rho$, pues a las entidades que esperan se suma la ocupación media del único servidor.

\subsection{Fórmula de Pollaczek--Khinchine e interpretación}
\label{subsec:pollaczek-khinchine}

Para una cola M/G/1 FCFS con los supuestos anteriores, la fórmula de Pollaczek--Khinchine establece
\begin{equation}
  \mathrm{E}[W_q]
  =\frac{\lambda\mathrm{E}[S^2]}{2(1-\rho)}
  =\frac{\rho m(1+C_s^2)}{2(1-\rho)}.
  \label{eq:pollaczek-khinchine}
\end{equation}
Virtamo desarrolla esta relación a partir del servicio residual y de la cantidad de entidades por delante de una llegada \cite{virtamoMG1}. Ambas expresiones tienen unidades de tiempo y describen una media estacionaria, no la espera de cada entidad.

La interpretación del servicio residual explica por qué interviene el segundo momento. Una llegada que encuentra ocupado al servidor debe esperar el resto de la atención en curso, además del servicio de quienes ya están en cola. Los servicios prolongados ocupan una mayor parte del tiempo y, por ello, tienen más probabilidad de ser encontrados por una observación temporal que los breves. No corresponde suponer que el servicio encontrado es una selección uniforme entre todas las atenciones. Las llegadas Poisson observan las proporciones temporales del estado del sistema, propiedad conocida como PASTA \cite{virtamoMG1}.

Si $R$ es el servicio residual observado, definido como cero cuando el servidor está libre, entonces
\begin{equation}
  \mathrm{E}[R]=\frac{\lambda\mathrm{E}[S^2]}{2},
  \qquad
  \mathrm{E}[W_q]=L_q m+\mathrm{E}[R].
  \label{eq:servicio-residual}
\end{equation}
Al sustituir Little y despejar se obtiene la ecuación~\ref{eq:pollaczek-khinchine}. Condicionado a encontrar ocupado al servidor, el residual medio es $\mathrm{E}[S^2]/(2m)$, no en general $m/2$. Esta diferencia recoge el mayor peso temporal de los servicios largos; incluso el servicio constante deja una espera residual cuando una llegada ocurre durante su ejecución.

Con $\lambda$ y $m$ fijos, también queda fija $\rho$. La espera media aumenta linealmente con la varianza del servicio, mientras que el factor $1/(1-\rho)$ amplifica la congestión al acercarse la utilización a uno. Dos distribuciones con iguales primeros dos momentos producen la misma espera media dentro de este modelo \cite{virtamoMG1}; no se deduce que tengan iguales distribuciones de espera, percentiles o probabilidades de demoras extremas.

\subsection{Casos particulares y alcance de la comparación}
\label{subsec:comparacion-casos}

El servicio determinista tiene $S=m$ y varianza nula; el exponencial de media $m$ tiene varianza $m^2$. Su sustitución en Pollaczek--Khinchine proporciona la comparación simbólica de la tabla~\ref{tab:casos-servicio}, manteniendo iguales $\lambda$ y $m$ \cite{virtamoMG1}.

\begin{table}[htbp]
  \centering
  \caption{Espera media teórica para dos casos de M/G/1 FCFS, con igual tasa de llegada y media de servicio.}
  \label{tab:casos-servicio}
  \begin{tabular}{lccc}
    \hline
    Modelo & Servicio & $C_s^2$ & $\mathrm{E}[W_q]$ \\
    \hline
    M/D/1 & Constante & $0$ & $\displaystyle\frac{\rho m}{2(1-\rho)}$ \\[6pt]
    M/M/1 & Exponencial & $1$ & $\displaystyle\frac{\rho m}{1-\rho}$ \\[6pt]
    \hline
  \end{tabular}
\end{table}

Bajo estas condiciones, la espera media con servicio exponencial es el doble de la correspondiente al servicio constante. La relación no se traslada directamente al tiempo total, porque ambos casos añaden la misma media $m$ de servicio propio. Además, como la varianza no puede ser negativa, el caso determinista minimiza la espera media entre las distribuciones admisibles de igual media. Son consecuencias analíticas del modelo, no hallazgos experimentales de esta monografía.

Este marco fija el lenguaje de la comparación posterior. La demostración prevista en Python, mediante un notebook compatible con Google Colab, deberá distinguir las identidades teóricas de las estimaciones obtenidas en ejecuciones finitas y considerar su incertidumbre. Aquí no se presentan simulaciones ejecutadas ni validación de un servicio real. La elección de una tercera distribución y el diseño de las réplicas corresponden a la metodología; las conclusiones sobre otras disciplinas o métricas exigirán evidencia específica.

% !TeX root = ../main.tex

\section{Revisión de la literatura}
\label{sec:revision-literatura}

\subsection{Una relación establecida y sus condiciones}
\label{subsec:literatura-relacion}

La bibliografía seleccionada permite distinguir dos preguntas: cómo influye la dispersión del servicio bajo condiciones constantes y qué intervenciones pueden mejorar un sistema concreto. La primera tiene una respuesta analítica para M/G/1 FCFS; la segunda exige identificar qué cambia realmente en el proceso. Esta revisión acotada organiza los aportes según esa distinción y sus límites, sin pretender una búsqueda sistemática ni una cobertura exhaustiva del campo.

La relación central no constituye una hipótesis teórica nueva. La ecuación~\ref{eq:pollaczek-khinchine}, desarrollada en el marco conceptual, establece que, a tasa de llegada y media de servicio fijas, una mayor varianza aumenta la espera media en cola, siempre que se cumplan los supuestos y la finitud de momentos indicados allí \cite{virtamoMG1}. El apéndice~1 de Joustra y colaboradores recupera esta relación para explicar el efecto de la variabilidad del servicio y la comparación entre servicio constante y exponencial \cite{joustra2012radioterapia}. Su utilidad para esta monografía es ofrecer una referencia común con la cual interpretar los demás trabajos, no una regla independiente de las condiciones del modelo.

El preprint de Lauber Bonomo, Pal y Reuveni también parte de M/G/1 FCFS y expresa el número medio de trabajos en el sistema mediante la utilización y el coeficiente de variación del servicio \cite{lauberbonomo2021resetting}. Es importante distinguir sus indicadores: en la sección~II, el tiempo denominado \emph{waiting time} comprende desde la llegada hasta el final del servicio. Corresponde, por tanto, a $T$ en nuestra notación, no a $W_q$. Con media de servicio fija, la diferencia entre ambos promedios es constante; cuando una intervención modifica esa media, la distinción vuelve a ser decisiva.

\subsection{Aplicaciones: variabilidad de capacidad y métodos complementarios}
\label{subsec:literatura-aplicaciones}

En procesos de varias etapas, hablar de variabilidad requiere precisar su origen. Joustra y colaboradores distinguen la dispersión de las duraciones individuales de servicio de las fluctuaciones de la capacidad disponible. En su estudio de radioterapia, la primera se relaciona con la espera en sala, mientras que la segunda afecta el acceso a la primera consulta y el flujo hacia preparación y tratamiento \cite{joustra2012radioterapia}. Esta separación evita interpretar cualquier reducción de demoras como evidencia de una disminución de la varianza del tiempo de atención.

El estudio combina análisis de datos del proceso, teoría de colas y simulación de eventos discretos. La teoría proporciona mecanismos explicativos; la simulación representa las interacciones entre etapas, localiza restricciones de capacidad y compara escenarios. Los autores contrastan el modelo con datos de la práctica para las medias y los percentiles considerados, y emplean intervalos de confianza al evaluar el cumplimiento de metas. En los escenarios estudiados, una capacidad de primeras consultas más estable permite mejorar el acceso y las demoras posteriores \cite{joustra2012radioterapia}. La evidencia cuantitativa procede del modelo de simulación contrastado con el proceso, no de un ensayo operativo que aísle únicamente la dispersión del servicio.

El límite de transferencia es concreto: se trata de un proceso con citas, recursos múltiples y etapas interdependientes, y varios escenarios incluyen un acelerador adicional. El propio artículo aclara que no interviene sobre la duración de las consultas, cuyas fluctuaciones eran pequeñas. Así, su aporte no consiste en comprobar experimentalmente nuestra comparación M/G/1 a media fija, sino en mostrar por qué conviene combinar explicación analítica y representación dinámica cuando la organización del servicio excede una cola aislada \cite{joustra2012radioterapia}. Para la demostración futura de esta monografía, la combinación tendrá un propósito más limitado: ilustrar una identidad conocida con parámetros hipotéticos, sin afirmar validación de un servicio real.

\subsection{Intervenciones que modifican el servicio efectivo}
\label{subsec:literatura-reinicios}

Otra forma de actuar sobre la congestión consiste en modificar el mecanismo de servicio. Lauber Bonomo y colaboradores analizan reinicios cuando la aleatoriedad es intrínseca al servidor: al reiniciar un trabajo, se obtiene una nueva duración aleatoria independiente, aunque el tiempo ya invertido se pierde. No equivale a adelantar otro cliente ni a eliminar trabajos pendientes. Si, por el contrario, cada trabajo conserva un requerimiento fijo y el reinicio solo obliga a repetirlo, el mecanismo puede desperdiciar tiempo sin aportar esa posibilidad de renovación \cite{lauberbonomo2021resetting}.

La estrategia matemática consiste en caracterizar el tiempo total de servicio efectivo, incluidos los intentos fallidos, y volver a aplicar el análisis M/G/1 con sus nuevos momentos. El preprint estudia reinicios a tasa constante y a intervalos fijos, y desarrolla ejemplos lognormales y de Pareto mediante fórmulas y evaluación numérica. Bajo las condiciones examinadas, la elección adecuada del reinicio puede reducir la media del servicio y su variabilidad relativa, con mejoras en el número medio de trabajos \cite{lauberbonomo2021resetting}. Son resultados analíticos y ejemplos numéricos del trabajo citado, no simulaciones realizadas en esta monografía.

Esta intervención no mantiene necesariamente constante $m$: también puede alterar $\rho=\lambda m$. En consecuencia, una mejora no debe atribuirse exclusivamente a la reducción de la varianza. Tampoco contradice que el servicio determinista minimice la espera media entre distribuciones de igual media: después del reinicio, la media efectiva puede ser diferente. El aporte del estudio es mostrar que dispersión y velocidad media pueden cambiar conjuntamente; nuestra comparación las separará deliberadamente \cite{lauberbonomo2021resetting}.

\subsection{Límites al cambiar estructura, disciplina o indicador}
\label{subsec:literatura-limites}

La estructura de la red también restringe las generalizaciones. En el resumen y el comienzo de la introducción de Bondi y Whitt se presentan principios sobre variabilidad para una red cerrada, de una clase de trabajos, con estaciones de un servidor y disciplina FCFS. Allí se anuncia que la longitud media en una estación que constituye un cuello de botella puede disminuir al aumentar la variabilidad del servicio en esa estación \cite{bondiwhitt1986variabilidad}. Esta referencia se utiliza únicamente como contraste de alcance, a partir de esa primera página consultada; no se evalúan aquí sus demostraciones, ejemplos numéricos ni procedimientos de aproximación. Una red cerrada no conserva el proceso externo de llegadas Poisson de nuestra cola abierta aislada, por lo que el principio anunciado no puede trasladarse directamente a ella.

El cambio de disciplina y de indicador ofrece un contraste más explícito. Koole, Nuyens y Righter estudian el máximo número de clientes en el sistema durante un período ocupado en M$^X$/G/1, con llegadas Poisson por lotes y variantes LCFS con interrupción: se atiende primero a la llegada más reciente. Sus resultados comparan distribuciones de ese máximo, no esperas medias FCFS. Para sistemas estables, igual tasa de llegada e igual distribución de tamaño del lote, obtienen máximos estocásticamente menores cuando los servicios están ordenados como más variables en sentido convexo \cite{koole2005maximos}.

El orden convexo exige comparar las esperanzas de todas las funciones convexas para las cuales existan; implica igualdad de medias y, cuando las varianzas son finitas, ordena también sus varianzas. Sin embargo, una varianza mayor por sí sola no establece ese orden. Por ello, los teoremas no autorizan a predecir el máximo a partir de cualquier comparación de varianzas. El manuscrito aporta caracterizaciones recursivas, demostraciones para las disciplinas preemptivas estudiadas e ilustraciones numéricas; reserva otras disciplinas para conjeturas \cite{koole2005maximos}. No corresponde presentar estas últimas como resultados demostrados ni interpretar un máximo menor como una mejora universal de todos los indicadores.

\subsection{Síntesis y fundamento de la comparación controlada}
\label{subsec:literatura-sintesis}

En conjunto, los trabajos no ofrecen conclusiones incompatibles: responden preguntas diferentes. La relación de referencia concierne a la espera media de una cola abierta FCFS; la aplicación hospitalaria modifica disponibilidad de capacidad; los reinicios transforman el servicio efectivo; los otros contrastes cambian la estructura de red o la disciplina y el indicador. Su lectura conjunta muestra que una afirmación sobre variabilidad debe identificar siempre qué se mantiene fijo, qué mecanismo se modifica y qué desempeño se mide.

La monografía se justifica, entonces, como una síntesis explicativa y una comparación controlada, no como la resolución de un vacío teórico. Mantener $\lambda$, $m$ y FCFS constantes permitirá aislar el efecto del segundo momento sobre $\mathrm{E}[W_q]$. La futura demostración deberá contrastar estimaciones con la referencia analítica y considerar su incertidumbre, sin equiparar una ejecución finita con la media estacionaria. La metodología precisará ese diseño; aquí queda fundamentado por qué conviene preservar estos controles y no extender sus conclusiones a redes, reinicios, máximos de cola o disciplinas no estudiadas.

% !TeX root = ../main.tex

\section{Metodología}
\label{sec:metodologia}

\subsection{Enfoque metodológico}
\label{subsec:metodo-enfoque}

La monografía combina una revisión narrativa acotada, una comparación analítica y el diseño de una demostración computacional complementaria. La selección y lectura de fuentes ya realizadas fundamentan las secciones anteriores; la programación, ejecución y evaluación del experimento se realizarán posteriormente. No se dispone de datos observados de una organización: los parámetros serán hipotéticos y la simulación no constituirá una validación de un servicio real.

\subsection{Escenarios y controles de comparación}
\label{subsec:metodo-escenarios}

Se representará una cola abierta aislada, con un servidor FCFS, capacidad de espera ilimitada, sin abandonos, prioridades ni interrupciones. Las llegadas serán Poisson de tasa constante $\lambda=0{,}8$ entidades por minuto. Los servicios serán independientes e idénticamente distribuidos e independientes de las llegadas. Se fijará $m=1$ minuto por entidad, de modo que $\rho=\lambda m=0{,}8<1$ en todos los escenarios. Estos controles permitirán concentrar el análisis de la espera media en el segundo momento del servicio, conforme a Pollaczek--Khinchine \cite{virtamoMG1}.

Se compararán servicio constante, exponencial e hiperexponencial de dos componentes, H$_2$. Para este último, cada entidad elegirá con probabilidad $0{,}8$ una exponencial de \emph{media} $m/2$ y con probabilidad $0{,}2$ una de \emph{media} $3m$. Estas cantidades no son tasas. La mezcla es una elección ilustrativa propia, no una ley ajustada a datos ni tomada de los artículos. La tabla~\ref{tab:metodo-servicios} resume parámetros poblacionales, no resultados simulados.

\begin{table}[htbp]
  \centering
  \caption{Distribuciones propuestas: parámetros y momentos teóricos del servicio.}
  \label{tab:metodo-servicios}
  \begin{tabular}{lccc}
    \hline
    Modelo y servicio & $\mathrm{E}[S]$ & $\operatorname{Var}(S)$ & $C_s^2$ \\
    \hline
    M/D/1: constante $m$ & $m$ & $0$ & $0$ \\
    M/M/1: exponencial, media $m$ & $m$ & $m^2$ & $1$ \\
    M/H$_2$/1: mezcla indicada & $m$ & $3m^2$ & $3$ \\
    \hline
  \end{tabular}
\end{table}

En la mezcla, $\mathrm{E}[S]=0{,}8(m/2)+0{,}2(3m)=m$ y $\mathrm{E}[S^2]=2[0{,}8(m/2)^2+0{,}2(3m)^2]=4m^2$; por tanto, su varianza es $3m^2$. Los tres segundos momentos son finitos. El tercer escenario será específicamente M/H$_2$/1, dentro de M/G/1, y no una distribución llamada G. La ecuación~\ref{eq:pollaczek-khinchine} proporciona referencias \emph{analíticas} de $2$, $4$ y $8$ minutos de espera media, respectivamente. No son mediciones ni resultados de una ejecución.

\subsection{Representación computacional y entorno previsto}
\label{subsec:metodo-computacion}

Se preparará un único notebook editable \texttt{.ipynb} para Google Colab, con explicaciones en celdas de texto, parámetros, código y análisis. NumPy se utilizará para generación pseudoaleatoria y arreglos; pandas, para tablas; Matplotlib, para gráficos. Se empleará además \texttt{scipy.stats.t} únicamente para obtener el cuantil crítico de Student. El uso previsto en Colab no requerirá configuración local. Se registrarán las versiones de Python y de los paquetes al ejecutar; si no pudiera comprobarse la ejecución completa en Colab, se informará esa limitación sin confundirla con una prueba local.

El simulador propio será pequeño: avanzará cronológicamente por las llegadas y calculará los eventos de inicio y fin de servicio. Si $X_i$ es un intervalo entre llegadas exponencial de media $1/\lambda$, se definirán
\begin{equation}
  A_i=\sum_{j=1}^{i}X_j,\qquad
  B_i=\max(A_i,D_{i-1}),\qquad D_i=B_i+S_i,
  \label{eq:metodo-eventos}
\end{equation}
con $D_0=0$, sistema inicialmente vacío. Para cada entidad se obtendrán $W_{q,i}=B_i-A_i$ y $T_i=D_i-A_i$. Esta recurrencia es equivalente a seguir los eventos de la cola de un servidor FCFS adoptada; no pretende reemplazar un calendario general para prioridades o redes. Se seguirá cada llegada seleccionada hasta su propia finalización, sin censurar servicios por un horizonte temporal fijo.

La figura~\ref{fig:procedimiento-simulacion} sintetiza el procedimiento previsto para el notebook. La evolución completa de cada realización se calculará antes de seleccionar las observaciones de las dos ventanas; el descarte afectará la estimación, no el estado de la cola. Las medias apareadas se resumirán entre réplicas de cada escenario para contrastar la teoría y examinar la sensibilidad al período inicial.

\begin{figure}[htbp]
  \centering
  % Fragmento editable: solo TikZ básico, sin bibliotecas adicionales.
% Esquema prospectivo; los bloques son etapas, no módulos implementados.
% Ancho: 0.82\textwidth más dos márgenes internos de 5 pt (<16 cm).
\begin{tikzpicture}[
  x=\textwidth,y=1cm,font=\small,
  etapa/.style={draw=black!75,fill=white,line width=0.5pt,
    text width=0.82\textwidth,align=center,inner sep=5pt,
    minimum height=0.95cm},
  flujo/.style={->,draw=black!75,line width=0.6pt}
]
  \node[etapa] (configuracion) at (0,0)
    {Configurar escenarios y semillas\\
     $\lambda=0{,}8$, $m=1$, $\rho=0{,}8$; constante, exponencial y H$_2$};
  \node[etapa] (generacion) at (0,-1.5)
    {Generar llegadas y servicios\\
     $140\,000$ entidades por réplica; flujos separados};
  \node[etapa] (simulacion) at (0,-3)
    {Simular FCFS desde vacío\\
     Calcular toda la realización: $A_i$, $B_i$, $D_i$, $W_{q,i}$};
  \node[etapa] (medicion) at (0,-4.5)
    {Descartar observaciones y medir sin reiniciar la cola\\
     Ventanas apareadas: medias $Y_r$ y $Y'_r$};
  \node[etapa] (resumen) at (0,-6)
    {Resumir $K=40$ réplicas por escenario\\
     Medias, intervalos y diferencias $\Delta_r=Y'_r-Y_r$};
  \node[etapa] (contraste) at (0,-7.5)
    {Comparar con Pollaczek--Khinchine\\
     Evaluar precisión y sensibilidad al período inicial};
  \draw[flujo] (configuracion.south) -- (generacion.north);
  \draw[flujo] (generacion.south) -- (simulacion.north);
  \draw[flujo] (simulacion.south) -- (medicion.north);
  \draw[flujo] (medicion.south) -- (resumen.north);
  \draw[flujo] (resumen.south) -- (contraste.north);
\end{tikzpicture}

  \caption{Esquema del procedimiento de simulación previsto, desde la configuración hasta el contraste analítico y el diagnóstico de inicialización. $Y_r$ y $Y'_r$ representan las medias de las ventanas principal y alternativa de una misma réplica. Elaboración propia: diseño previsto, no resultados ni software ya ejecutado.}
  \label{fig:procedimiento-simulacion}
\end{figure}

\subsection{Réplicas y reproducibilidad pseudoaleatoria}
\label{subsec:metodo-replicas}

Se realizarán $K=40$ réplicas independientes por escenario. Cada una generará $140\,000$ llegadas y servicios desde vacío. La estimación principal descartará las primeras $b=20\,000$ llegadas y utilizará las siguientes $n=100\,000$, índices $b+1,\ldots,b+n$. Descartar observaciones no reiniciará la cola: se conservarán el trabajo pendiente y la evolución acumulada. El sufijo adicional permitirá el diagnóstico de inicialización descrito más adelante. Se fija un presupuesto por cantidad de entidades, no por minutos simulados; estos valores son decisiones iniciales de diseño y no demuestran estacionariedad ni precisión suficiente.

La semilla maestra será 2026. Para cada escenario y réplica se construirá una \texttt{SeedSequence} de NumPy con los enteros \texttt{[2026, escenario, replica]}, donde los escenarios tendrán identificadores $0,1,2$ y las réplicas $0,\ldots,39$. Desde ella se derivarán mediante \texttt{spawn} flujos separados para intervalos de llegada, elección del componente y duraciones de servicio. Se mantendrá esta asignación incluso cuando un escenario no necesite elegir componentes. Se asumirá independencia práctica entre flujos pseudoaleatorios, no una garantía criptográfica. No se usarán números aleatorios comunes entre escenarios; el apareamiento se reservará al diagnóstico dentro de una misma realización.

\subsection{Estimación e incertidumbre estadística}
\label{subsec:metodo-inferencia}

La unidad del análisis estadístico será la media de cada réplica,
\begin{equation}
  Y_r=\frac{1}{n}\sum_{i=b+1}^{b+n}W_{q,i}^{(r)},\qquad
  \overline{Y}=\frac{1}{K}\sum_{r=1}^{K}Y_r,\qquad
  s_Y^2=\frac{\sum_{r=1}^{K}(Y_r-\overline{Y})^2}{K-1}.
  \label{eq:metodo-estimadores}
\end{equation}
Las esperas sucesivas dentro de una réplica pueden estar correlacionadas: no se tratarán como observaciones independientes para calcular el error estándar. El método de réplicas permite estimar la variabilidad entre medias, distinguiéndola de la dependencia interna \cite{banksAnalisisSalidas}. Se construirá el intervalo aproximado del $95\,\%$
\begin{equation}
  \overline{Y}\pm h,\qquad
  h=t_{0{,}975;K-1}\frac{s_Y}{\sqrt{K}},
  \label{eq:metodo-intervalo}
\end{equation}
donde $t_{0{,}975;K-1}$ es el cuantil de Student con $K-1$ grados de libertad. Su aproximación dependerá de la distribución de las medias de réplica. El intervalo describirá variabilidad Monte Carlo condicionada al modelo y al diseño, no sesgo residual, validez del modelo ni una cobertura exacta universal.

Se informará la semiamplitud relativa $h/\overline{Y}$, con meta diagnóstica de no superar el $5\,\%$. Si el presupuesto fijado no alcanzara esa precisión, se declarará la limitación, sin prolongar automáticamente las ejecuciones. También se calculará el error relativo $(\overline{Y}-w_\star)/w_\star$, donde $w_\star$ es la referencia analítica correspondiente. La inclusión de esa referencia en el intervalo será descriptiva: ni error exactamente nulo ni cobertura en toda ejecución serán requisitos de aprobación. La comparación principal será analítica; las simulaciones la ilustrarán, sin inferir diferencias entre modelos a partir del mero solapamiento de intervalos marginales.

\subsection{Sensibilidad al período inicial}
\label{subsec:metodo-inicializacion}

El inicio vacío puede introducir sesgo frente al régimen estacionario; aumentar las réplicas no lo elimina por sí mismo \cite{banksAnalisisSalidas}. Se evaluará una alternativa con descarte $b'=40\,000$ y las siguientes $100\,000$ llegadas de la \emph{misma} realización. Para cada réplica se obtendrá $Y'_r$ y la diferencia apareada $\Delta_r=Y'_r-Y_r$. Las ventanas se superpondrán, por lo que no serán independientes entre sí; las diferencias sí se analizarán como independientes entre réplicas.

La figura~\ref{fig:ventanas-medicion} sitúa ambas ventanas en un único eje de índices de entidad. El tramo común de medición comprende las entidades $40\,001$ a $120\,000$; las últimas $20\,000$ solo intervienen en la media alternativa. No habrá una segunda simulación ni un reinicio de estado al cambiar de ventana.

\begin{figure}[htbp]
  \centering
  % Fragmento editable: solo TikZ básico, sin bibliotecas adicionales.
% Escala común proporcional: x = índice / 140000.
% Límites exactos: 0, 1/7, 2/7, 6/7, 1.
% Las barras ocupan 0.84\textwidth; los extremos dejan margen para rótulos.
% Los descartes se rotulan fuera de los segmentos cortos, sin reducir letra.
\begin{tikzpicture}[
  x=0.84\textwidth,y=1cm,font=\small,
  borde/.style={draw=black!75,line width=0.5pt},
  descarte/.style={borde,fill=black!8,dashed},
  medicion/.style={borde,fill=blue!10},
  nota/.style={align=center,inner sep=2pt}
]
  \node[nota] at (0.5,2.5)
    {Una realización por réplica; evolución continua de la cola};
  \node[anchor=west,inner sep=0pt] at (0,1.95) {Ventana principal: $Y_r$};
  \path[descarte] (0,1) rectangle ({1/7},1.6);
  \path[medicion] ({1/7},1) rectangle ({6/7},1.6);
  \path[borde,fill=white] ({6/7},1) rectangle (1,1.6);
  \node[nota] at (0.5,1.3) {Medición: $100\,000$};
  \draw[borde] ({1/14},1) -- ({1/14},0.8);
  \node[nota,anchor=north] at ({1/14},0.8) {Descarte\\$20\,000$};
  \draw[borde] ({13/14},1) -- ({13/14},0.8);
  \node[nota,anchor=north] at ({13/14},0.8) {No usado\\en $Y_r$};

  \node[anchor=west,inner sep=0pt] at (0,-0.15) {Ventana alternativa: $Y'_r$};
  \path[descarte] (0,-1.1) rectangle ({2/7},-0.5);
  \path[medicion] ({2/7},-1.1) rectangle (1,-0.5);
  \node[nota] at ({9/14},-0.8) {Medición: $100\,000$};
  \draw[borde] ({1/7},-1.1) -- ({1/7},-1.3);
  \node[nota,anchor=north] at ({1/7},-1.3) {Descarte\\$40\,000$};

  % Proyecciones de los límites sobre un único eje de índices, no temporal.
  \foreach \x in {0,{1/7},{2/7},{6/7},1} {
    \draw[draw=black!40,line width=0.45pt,dotted]
      (\x,-1.95) -- (\x,-2.2);
  }
  \draw[borde,->] (0,-2.2) -- (1.025,-2.2);
  \foreach \x/\valor in {0/0,{1/7}/{20\,000},{2/7}/{40\,000},
      {6/7}/{120\,000},1/{140\,000}} {
    \draw[borde] (\x,-2.15) -- (\x,-2.25);
    \node[anchor=north,inner sep=2pt] at (\x,-2.28) {$\valor$};
  }
  \node[nota,anchor=north] at (0.5,-2.85) {Índice de entidad};
\end{tikzpicture}

  \caption{Esquema de las ventanas apareadas sobre una misma realización de $140\,000$ entidades. Cada bloque delimitado por $a$ y $b$ comprende los índices $a+1,\ldots,b$; el eje cuenta entidades, no minutos. El descarte principal abarca $1$--$20\,000$ y la medición $20\,001$--$120\,000$; la alternativa descarta $1$--$40\,000$ y mide $40\,001$--$140\,000$. La cola evoluciona sin reinicios. Elaboración propia: diseño previsto y esquema del procedimiento, no resultados.}
  \label{fig:ventanas-medicion}
\end{figure}

Se calcularán $\overline{\Delta}$, su desviación muestral $s_\Delta$ y el intervalo aproximado $\overline{\Delta}\pm h_\Delta$, con $h_\Delta=t_{0{,}975;K-1}s_\Delta/\sqrt{K}$. Se usará el criterio práctico
\begin{equation}
  |\overline{\Delta}|+h_\Delta\leq 0{,}02\,w_\star.
  \label{eq:metodo-sensibilidad}
\end{equation}
Esta comprobación acotará la diferencia detectable respecto de un umbral propio del diseño; no probará ausencia de sesgo, pues ambas ventanas podrían compartirlo. Si no se satisficiera, se informará incertidumbre residual de inicialización, sin reintentos ilimitados ni declaración automática de convergencia. Los tamaños y umbrales propuestos no se atribuyen como recomendaciones numéricas de las diapositivas citadas.

\subsection{Verificación prevista y productos trazables}
\label{subsec:metodo-verificacion}

Antes de interpretar salidas, se incorporarán pruebas deterministas visibles y ejecutables desde el notebook. Con llegadas $A=[0;0{,}5;2]$ y servicios $S=[1;1;1]$, deberán obtenerse inicios $B=[0;1;2]$, salidas $D=[1;2;3]$ y esperas $[0;0{,}5;0]$. El caso $A=[0;2;4]$, con iguales servicios, deberá producir esperas nulas. Se comprobarán esperas no negativas, orden FCFS, ausencia de servicios superpuestos y la identidad $T_i=W_{q,i}+S_i$. Las medias y varianzas muestrales de entrada se contrastarán con sus momentos conocidos como diagnóstico, no exigiendo igualdad exacta.

Después se contrastarán medias estimadas, intervalos y errores con Pollaczek--Khinchine. Esta verificación de consistencia con un modelo conocido no validará un servicio real. Se conservarán CSV con configuración y métricas por réplica, y una tabla resumen con media, varianza entre réplicas, intervalo, referencia, error y diagnósticos de precisión e inicialización. El notebook, las semillas y las versiones permitirán reconstruir el procedimiento.

Se prepararán cuatro figuras complementarias: las funciones de distribución acumulada teóricas de los servicios, la espera media simulada con sus intervalos y referencias analíticas, la distribución acumulada descriptiva de las esperas y la sensibilidad al descarte inicial. La primera será analítica, con un salto en $m$ para el servicio constante; las demás se basarán en salidas efectivamente obtenidas. Para describir las esperas se acumularán, en la ventana principal de cada réplica, conteos en $321$ umbrales entre $0$ y $100m$. Se dividirá su suma por $Kn$ y se informará la fracción superior al último umbral. La curva entre umbrales será una representación discretizada, no la distribución empírica completa; no se supondrá independencia entre esperas ni se añadirán intervalos por observación. Este procesamiento descriptivo no cambiará los flujos aleatorios ni los estimadores principales.

El notebook incorpora una conexión explícita a Google Drive que requiere la autorización del usuario en Colab. Se creará una carpeta distinta por ejecución, con configuración y estado iniciales; las métricas y los conteos se guardarán tras cada réplica. Los archivos incompletos se identificarán como parciales y el estado solo indicará finalización después de exportar los productos. Las cuatro figuras se guardarán como PNG junto con sus datos CSV; no se exportará a PDF sin autorización. La versión anterior ejecutada por el usuario se conserva por separado. La ejecución integral, la persistencia en Drive y la revisión visual de la versión ampliada quedan pendientes de comprobación; su implementación no constituye por sí sola evidencia de resultados.


\clearpage
\phantomsection
\addcontentsline{toc}{section}{\refname}
\bibliographystyle{unsrt}
\bibliography{referencias}

\end{document}
